% !TeX program = xelatex
% 在report目录编译两次：xelatex -interaction=nonstopmode -halt-on-error report.tex
% 图片使用现有outputs/figures文件；本轮仅改写tex，未运行回测或编译。
\documentclass[UTF8,a4paper,11pt,fontset=fandol]{ctexart}
\usepackage[margin=2.4cm]{geometry}
\usepackage{amsmath,amssymb,booktabs,tabularx,longtable,array}
\usepackage{graphicx,float,caption,enumitem,xcolor,xurl}
\usepackage[unicode,hidelinks]{hyperref}
\graphicspath{{../outputs/figures/}}
\setlength{\parindent}{2em}
\setlength{\parskip}{0.3em}
\renewcommand{\arraystretch}{1.18}
\setlist[itemize]{leftmargin=2em,itemsep=0.2em,topsep=0.3em}
\captionsetup{font=small,labelfont=bf}
\newcolumntype{Y}{>{\raggedright\arraybackslash}X}
\newcommand{\code}[1]{\texttt{\detokenize{#1}}}
\newcommand{\snapshotfigure}[3]{%
  \begin{figure}[htbp]\centering
  \includegraphics[width=\linewidth]{#1}
  \caption{#2}\label{#3}
  \end{figure}}
\newcommand{\pairfigure}[5]{%
  \begin{figure}[htbp]\centering
  \includegraphics[width=\linewidth]{#1}\par\smallskip
  \includegraphics[width=\linewidth]{#2}
  \caption{#3。上图：#4；下图：#5。}
  \end{figure}}
\title{基于平均真实波幅ATR的ETF网格交易策略\\收益、风险与择时效果分析}
\author{吴乙弘}
\date{\today}
\begin{document}
\maketitle
\begin{abstract}
ATR网格以价格偏离均线的波动幅度决定仓位：相对高位减仓、相对低位加仓，并以波动趋势和价格趋势调整交易方向。本文在六个A股宽基指数及其对应ETF上研究这一方法，指标窗口为20日，趋势判断窗口为60日，样本截至2025年9月30日。结果显示，基础网格显著降低波动与回撤，但在2019年以来的上涨收益积累中明显落后买入持有。加入择时后，沪深300的年化收益由2.89\%提高至最高5.00\%，最大回撤由23.24\%降至20.80\%；不过，改善幅度和风险效果随指数而异。固定六ETF组合采用中枢调整规则，年化收益5.52\%、最大回撤24.09\%、年化波动10.67\%，2022、2023年分别亏损13.31\%和2.20\%。相较广发金工研报宽基组合的11.78\%年化收益与11.50\%最大回撤，本研究仍存在明显差距。其投资含义是：网格可以改变权益风险暴露，但持续上涨中的收益让渡与下跌趋势中的仓位累积仍然存在，仅靠多ETF配置和简单择时尚不足以获得稳定的年度正收益。
\end{abstract}
\tableofcontents
\newpage

\section{研究思路与主要发现}
网格交易的基本逻辑是利用价格往复波动，逢低增加持仓、逢高降低持仓。固定价格间距在低波动时可能长期不触发交易，在高波动时又可能过早用尽仓位。平均真实波幅（Average True Range，ATR）为交易间距提供了随市场波动变化的尺度：相同的价格偏离，在高波动环境中被视为较小偏离，在低波动环境中则被视为较大偏离。

本文参考广发金工《基于平均真实波幅（ATR）的ETF网格交易策略》\cite{gfreport}，先在指数上研究仓位规则，再将指数信号映射到ETF，纳入分红和份额折算后构建六宽基组合。指数层面比较基础网格及四种独立择时方法；组合层面采用其中的中枢调整方法。研究围绕三个问题展开：降低回撤需要付出多大收益代价，择时能否改善这一取舍，以及多个宽基ETF组合后能否提高收益稳定性。

研究结果有三点。第一，基础网格在六个指数上均降低了2019年以来的最大回撤，也均落后于对应指数的年化收益。它较适合被理解为一种权益仓位管理方法，较低回撤本身不足以证明具有超额收益。第二，趋势约束可以缓解部分过早减仓和逆势加仓，但改善并不普遍；同一规则可能在沪深300上同时改善收益和回撤，在创业板或中证1000上却扩大回撤。第三，固定六ETF组合没有重现研报的年度正收益特征，差距主要体现在2022年的亏损和2024年的上涨参与不足；这些年度差异比单一的全样本年化收益更能揭示策略局限。

\section{数据与回测设计}
\subsection{研究资产与样本区间}
研究资产覆盖上证50、沪深300、中证500、中证1000、创业板指和科创50。指数日线使用开盘、最高、最低和收盘价，其中最高、最低及前收盘价用于计算ATR；ETF使用收盘价计算持有收益。项目数据下载接口分别为AKShare提供的腾讯指数历史日线和新浪ETF历史日线，分红数据采用新浪累计分红记录；回测使用本地日频数据。研报图表的数据源为Wind，两者未作逐日一致性匹配。

\begin{table}[htbp]\centering\small
\caption{指数与ETF对应关系及行情历史起点；数据截止日均为2025-09-30}
\begin{tabularx}{\linewidth}{l l l Y l}\toprule
指数 & 指数代码 & 指数历史起点 & 对应ETF & ETF行情起点\\\midrule
上证50 & 000016 & 2004-01-02 & 华夏上证50ETF（510050） & 2005-02-23\\
沪深300 & 000300 & 2005-04-08 & 华泰柏瑞沪深300ETF（510300） & 2012-05-28\\
中证500 & 000905 & 2007-01-15 & 南方中证500ETF（510500） & 2013-03-15\\
中证1000 & 000852 & 2014-10-17 & 南方中证1000ETF（512100） & 2016-11-04\\
创业板指 & 399006 & 2010-06-02 & 易方达创业板ETF（159915） & 2011-12-09\\
科创50 & 000688 & 2020-01-02 & 华夏科创50ETF（588000） & 2020-11-16\\\bottomrule
\end{tabularx}
\end{table}

主样本从2019年1月2日至2025年9月30日，另以2022年1月4日为起点观察后半段市场环境。指标使用样本开始前的历史数据预热。科创50历史较短，其指数策略在指标和滞后仓位均有效后，从2020年5月8日起统计；588000独立策略从2020年11月17日起统计。因此，科创50相关结果与其余五个标的的长期收益不宜直接按累计收益排序。2025年的年度收益均仅覆盖1--9月。

指数回测的基准为指数价格收益；ETF回测的基准为计入分红和份额折算的ETF买入持有收益。前者用于比较配置规则，后者更接近投资工具的持有回报，二者不能互相替代。

\subsection{分红、份额折算与总回报}
ETF分红使用累计分红相邻两次记录之差还原单次现金分红，并按除息日计入持有收益。例如，510300在2017年1月23日的单次分红为$0.222-0.167=0.055$元/份。510050、510300、510500和512100使用各自分红记录；159915和588000按研究区间无分红处理。分红视为除息日立即进入总回报，未进一步模拟实际到账日及再投资成交。

份额合并或拆分会造成单位价格跳变，必须按持有人总资产价值计算收益。设$P_{i,t}$为ETF收盘价，$q_{i,t}$为每份旧份额换得的新份额数（普通交易日为1），$D_{i,t}$为每份现金分红，则
\[
r_{i,t}^{ETF}=\frac{q_{i,t}P_{i,t}+D_{i,t}}{P_{i,t-1}}-1.
\]
\begin{table}[htbp]\centering\small
\caption{ETF份额折算及其对单日收益的影响}
\begin{tabularx}{\linewidth}{l l r r Y}\toprule
ETF & 调整日 & 份额换算系数 & 原始价格收益 & 调整后收益\\\midrule
510500 & 2015-04-15 & 0.2803 & --- & 早于业绩统计区间\\
510500 & 2022-08-29 & 1.1454 & -12.7424\% & -0.0551\%\\
512100 & 2022-09-05 & 0.36555 & 176.2729\% & 0.9916\%\\\bottomrule
\end{tabularx}
\end{table}
折算调整避免将份额变化误判为投资盈亏。调整后，2019年以来六ETF均无绝对值超过35\%的单日收益。不过，收益阈值只能识别明显异常，分红历史的完整性、除息日与行情日期的匹配，以及未登记的公司行动仍可能影响总回报。

\subsection{交易与绩效口径}
\begin{table}[htbp]\centering
\caption{主要回测假设}
\begin{tabularx}{\linewidth}{l l Y}\toprule
项目 & 设定 & 经济含义\\\midrule
ATR与价格中枢窗口 & 20个交易日 & 衡量近期价格偏离和日内、隔夜波动\\
趋势判断窗口 & 60个交易日 & 平滑波动指标与价格，识别方向\\
基础仓位/调整幅度 & 50\%/每档10个百分点 & 保留加减仓空间，仓位限于0--100\%\\
交易时点 & 信号次日收盘执行 & 新仓位从信号后第二个交易日起计收益\\
交易费用与滑点 & 0 & 结果未扣除实际交易摩擦\\
闲置资金收益 & 0 & 未配置货币或短融资产\\
年化换算 & 252个交易日 & 夏普无风险收益设为0\\\bottomrule
\end{tabularx}
\end{table}

年化收益按净值复合增长率换算，最大回撤为相对历史最高净值的最大跌幅，年化波动率为日收益样本标准差乘$\sqrt{252}$，夏普为年化算术平均收益除以年化波动率。以下回撤以正数展示，数值越小表示回撤越低。

指数和单ETF策略在样本开始时继承历史信号对应的仓位；六ETF组合从空仓开始，前两个交易日尚未形成计收益持仓。这一期初差异影响两类账户的可比性。组合还将缺失ETF行情日的收益记为0，未限制停牌期间的模拟调仓，因此属于按权重计算的历史模拟，尚未完整体现实际可成交条件。

\section{策略规则：从波动尺度到资产配置}
\subsection{基础网格：围绕移动中枢逆向调整仓位}
真实波幅同时考虑当日价格振幅和相对前收盘价的跳空。设$H_t,L_t,C_t$为最高、最低和收盘价，$MA$为简单移动平均，则
\begin{align}
TR_t&=\max\{H_t-L_t,|H_t-C_{t-1}|,|L_t-C_{t-1}|\},\\
ATR_t&=MA(TR_t,20),\qquad MC_t=MA(C_t,20),\\
NATR_t&=\frac{C_t-MC_t}{ATR_t}.
\end{align}
$NATR$表示价格高于或低于20日均线多少个ATR，正值表示相对高位，负值表示相对低位。这里的$NATR$是价格偏离指标，与常见的$ATR/C\times100\%$波动率指标不同。

当$-1<NATR<1$时持仓50\%；价格每向上偏离一档，减仓10个百分点，每向下偏离一档，加仓10个百分点。$NATR\geq5$时空仓，$NATR\leq-5$时满仓。例如$NATR=1.74$对应40\%仓位，$NATR=-3.2$对应80\%仓位。中枢随均线移动，交易间距随ATR变化，因此本策略的“网格”是每天更新的仓位映射，而非围绕固定价格挂出的永久买卖单。

这一规则在价格往复时能够反向调节风险，但也存在方向性代价：持续上涨时逐步降低参与度，持续下跌时逐步增加风险敞口。网格能否改善风险收益，取决于价格回归是否足以补偿趋势中的收益让渡或损失累积。

\subsection{ATR择时：用趋势约束逆向交易}
ATR择时先判断波动是否收敛，再判断价格方向。令$A_t=MA(ATR_t,60)$、$P_t=MA(C_t,60)$：当$A_t$下降且$P_t$上升时看多，当二者同时下降时看空，其余情形为中性。看多并不意味着直接满仓，看空也不意味着直接清仓，而是改变基础网格的调仓方式。沪深300主样本中，看多482天、看空458天、中性699天。

\begin{table}[htbp]\centering\small
\caption{四种独立择时规则及对应仓位变量}
\begin{tabularx}{\linewidth}{l Y Y l}\toprule
规则 & 看多时 & 看空时 & 仓位变量\\\midrule
策略1：调整间距 & 减仓间距扩大至2倍，加仓间距缩至0.5倍 & 减仓间距缩至0.5倍，加仓间距扩大至2倍 & \code{w_t1}\\
策略2：调整中枢 & 仓位中枢由50\%升至60\% & 仓位中枢降至40\% & \code{w_t2}\\
策略3：约束上下限 & 仓位不低于50\% & 仓位不高于50\% & \code{w_t3}\\
策略4：约束调仓方向 & 仓位不得低于前一期 & 仓位不得高于前一期 & \code{w_t4}\\\bottomrule
\end{tabularx}
\end{table}

中性状态下均恢复基础网格规则。策略1使交易对趋势有不同敏感度；策略2温和调整整体风险预算；策略3明确限制看多时过度减仓和看空时过度加仓；策略4则延续上一期仓位方向，具有路径依赖。四种规则分别回测，不相互叠加。基础策略对应\code{w_base}，具体代码位置见附录\ref{app:code}。

策略2沿用每侧最多五档的设计，因此看多时仓位范围为10--100\%，看空时为0--90\%。策略4使用自身上期仓位递推，而非上期基础仓位；2022年子样本继续继承此前状态。

\subsection{六宽基ETF组合：等额上限下的动态权益仓位}
组合采用策略2，将每个指数的目标仓位映射到对应ETF。若当日已有$K_t$只ETF上市且指数信号有效，则每只ETF的资金上限为$1/K_t$。设$e_{i,t}$为纳入标识，$w_{i,t}^{(2)}$为策略2仓位，则ETF目标权重为
\[
a_{i,t}=\frac{e_{i,t}w_{i,t}^{(2)}}{K_t},\qquad \sum_i a_{i,t}\leq1.
\]
这里“等额上限”不代表每只ETF始终等权持有：实际权重还随各自网格仓位变化，剩余部分保持现金。2019年有五只ETF参与，单只上限20\%；588000于2020年11月16日上市后，有效数量增至六只，单只上限降至16.67\%。

信号在$t$日收盘形成，$t+1$日收盘调仓，从$t+2$日起按新权重计入ETF总回报。组合不依据基金规模和成交额动态更换ETF，也不在六类宽基之外寻找资产。因此，它检验的是固定宽基工具下的网格配置效果，未包含全市场ETF择优与行业轮动。

\clearpage
\section{指数结果：降低回撤伴随上涨收益让渡}
\subsection{沪深300长期样本：择时改善了基础网格的取舍}
2019年以来，沪深300基础网格累计收益20.33\%，低于指数的54.14\%；同时，最大回撤从指数的45.60\%降至23.24\%，年化波动由19.22\%降至9.58\%。收益和风险同时下降，与较低权益暴露的机制一致。基础网格夏普为0.3449，低于指数的0.4424，说明该样本中风险下降并未带来更高的单位波动收益。

\begin{table}[htbp]\centering\small
\caption{沪深3002019-01-02至2025-09-30的策略表现}
\begin{tabular}{lrrrrr}\toprule
方法 & 累计收益 & 年化收益 & 最大回撤 & 年化波动 & 夏普\\\midrule
基础网格 & 20.33\% & 2.89\% & 23.24\% & 9.58\% & 0.3449\\
策略1：间距 & 33.16\% & 4.50\% & 20.41\% & 9.57\% & 0.5080\\
策略2：中枢 & 26.55\% & 3.69\% & 21.54\% & 9.55\% & 0.4270\\
策略3：上下限 & 37.37\% & 5.00\% & 20.80\% & 9.66\% & 0.5537\\
策略4：方向 & 34.10\% & 4.61\% & 21.28\% & 9.64\% & 0.5162\\
指数买入持有 & 54.14\% & 6.88\% & 45.60\% & 19.22\% & 0.4424\\\bottomrule
\end{tabular}
\end{table}

四种择时均提高了基础网格收益，并降低最大回撤。其中，策略3年化收益最高，较基础提高约2.12个百分点；策略1最大回撤最低，为20.41\%。策略3在波动率仅略高于基础的情况下，将夏普提高至0.5537，也高于指数。这表明，在沪深300这一特定样本中，保留看多时的最低参与度、约束看空时的最高仓位，比纯粹逆向加减仓取得了更好的收益与风险组合。

但策略3年化收益仍比指数低约1.88个百分点，累计收益少约16.77个百分点。择时缓解了上涨中的收益损失，没有消除网格对持续趋势的参与不足。此外，缺少固定50\%指数仓位这一对照，尚不能将全部风险改善归因为网格交易本身。
\pairfigure{index_000300_2019_base_vs_index.png}{index_000300_2019_five_strategies.png}{沪深3002019年以来策略净值}{基础网格与指数}{基础网格及四种择时规则}

\subsection{沪深3002022年以来：防守改善仍未转为正收益}
2022年以来，基础网格累计收益为$-6.57\%$，略低于指数的$-6.07\%$，但最大回撤低约15.92个百分点。四种择时均优于基础及指数，策略4的累计损失最小，为$-2.99\%$，最大回撤也最低，为16.07\%。
\begin{table}[htbp]\centering\small
\caption{沪深3002022-01-04至2025-09-30的策略表现}
\begin{tabular}{lrrrrr}\toprule
方法 & 累计收益 & 年化收益 & 最大回撤 & 年化波动 & 夏普\\\midrule
基础网格 & -6.57\% & -1.87\% & 20.13\% & 9.21\% & -0.1585\\
策略1：间距 & -3.43\% & -0.96\% & 17.42\% & 8.99\% & -0.0625\\
策略2：中枢 & -3.11\% & -0.87\% & 17.14\% & 8.99\% & -0.0524\\
策略3：上下限 & -3.02\% & -0.85\% & 16.88\% & 8.89\% & -0.0513\\
策略4：方向 & -2.99\% & -0.84\% & 16.07\% & 8.91\% & -0.0499\\
指数买入持有 & -6.07\% & -1.72\% & 36.05\% & 18.14\% & -0.0052\\\bottomrule
\end{tabular}
\end{table}

这一结果说明，趋势约束在该区间能够减少持续下跌中的损失，但保护是相对的：所有策略期末仍未回到期初资产水平。长期样本以策略3收益最好，后半段样本则以策略4最好，也表明规则排序会随市场路径变化，不能根据单一区间确定稳定最优方案。
\pairfigure{index_000300_2022_base_vs_index.png}{index_000300_2022_five_strategies.png}{沪深3002022年以来策略净值}{基础网格与指数}{基础网格及四种择时规则}

\subsection{创业板与其他指数：收益改善并不总伴随回撤下降}
创业板的结果更突出地体现了网格在趋势市场中的代价。2019年以来，指数累计上涨158.94\%，基础网格仅上涨36.69\%；策略3虽将累计收益提高至61.98\%，仍明显落后指数。其夏普0.6032也低于指数的0.6461，与沪深300中策略3夏普超过指数的结果不同。
\begin{table}[htbp]\centering\small
\caption{创业板指两个区间的代表性结果}
\begin{tabular}{llrrrr}\toprule
区间 & 方法 & 累计收益 & 年化收益 & 最大回撤 & 夏普\\\midrule
2019年以来 & 基础网格 & 36.69\% & 4.92\% & 34.50\% & 0.4124\\
 & 策略3 & 61.98\% & 7.70\% & 34.27\% & 0.6032\\
 & 指数 & 158.94\% & 15.75\% & 57.05\% & 0.6461\\
2022年以来 & 基础网格 & -11.92\% & -3.46\% & 32.45\% & -0.1755\\
 & 策略3 & -8.01\% & -2.29\% & 32.27\% & -0.1029\\
 & 指数 & -2.54\% & -0.71\% & 53.94\% & 0.1228\\\bottomrule
\end{tabular}
\end{table}

创业板2019年以来的四种择时均提高了基础收益，但策略1、2、4的最大回撤分别增至35.61\%、35.48\%和35.29\%，高于基础的34.50\%。趋势信号增加部分时期的参与度，也可能增加后续损失；“加入择时”并不自动意味着收益与风险同时改善。2022年以来，四种择时均未在期末跑赢指数；即使策略3累计亏损已较基础缩小，仍比指数多亏约5.46个百分点。

六指数的横向比较进一步表明，规则效果存在明显差异。中证1000在2019年以来采用策略1、2、4后，年化收益由基础的4.54\%分别降至4.20\%、4.08\%、4.08\%，最大回撤均扩大。策略3在六个长期样本中都提高了基础年化收益，但科创50的回撤由36.43\%升至37.04\%。科创50长期样本从2020年开始，相关比较应在自身相同起止日内进行。

因此，跨指数结果支持的是“部分趋势约束可以改善部分网格策略”，而非任何择时规则均有效。策略3在本组长期样本中较有一致性，但仍需区分样本内表现与未来可持续性，也不能据此推定六ETF采用策略3后会获得何种组合收益。
\pairfigure{index_399006_2019_base_vs_index.png}{index_399006_2019_five_strategies.png}{创业板指2019年以来策略净值}{基础网格与指数}{基础网格及四种择时规则}
\pairfigure{index_399006_2022_base_vs_index.png}{index_399006_2022_five_strategies.png}{创业板指2022年以来策略净值}{基础网格与指数}{基础网格及四种择时规则}

\clearpage
\section{ETF结果：趋势参与不足在可交易工具上仍然存在}
将相同指数信号映射到ETF并计入分红和份额折算后，前五只ETF均以策略3获得最高的网格累计收益；但全部六只ETF的网格策略均落后于各自的买入持有基准。网格与基准的差距并未因总回报口径而消失。
\begin{table}[htbp]\centering\small
\caption{六ETF独立策略累计收益；前五只从2019年起，588000从2020-11-17起}
\begin{tabular}{lrrrrrr}\toprule
ETF & 基础网格 & 策略1 & 策略2 & 策略3 & 策略4 & 买入持有\\\midrule
510050 & 26.69\% & 30.84\% & 32.62\% & 34.09\% & 33.89\% & 50.19\%\\
510300 & 28.90\% & 42.39\% & 35.15\% & 46.87\% & 42.74\% & 75.01\%\\
510500 & 47.49\% & 47.78\% & 46.74\% & 53.26\% & 47.39\% & 99.78\%\\
512100 & 51.40\% & 46.03\% & 45.89\% & 59.74\% & 45.60\% & 115.05\%\\
159915 & 38.74\% & 48.63\% & 40.45\% & 65.64\% & 51.06\% & 169.43\%\\
588000 & 1.70\% & 3.53\% & 3.60\% & 0.73\% & -2.73\% & 7.17\%\\\bottomrule
\end{tabular}
\end{table}

以510300为例，策略3累计收益46.87\%，较基础的28.90\%高约17.97个百分点，但仍比买入持有少约28.14个百分点。159915的差距更大，策略3收益65.64\%，而买入持有为169.43\%。这一现象与指数结果一致：较低权益仓位改善持有路径的同时，也会让渡较多长期上涨收益。

中证1000ETF的策略1、2、4均低于基础，说明调整交易节奏或中枢没有形成普遍收益增益。科创50ETF则以策略2累计收益3.60\%最高，策略3仅0.73\%，策略4亏损2.73\%；“策略3在其他标的较好”的结论无法直接推广至不同上市时点和价格路径的资产。

ETF基准包含分红，指数基准是价格指数，因此两者终值不同。例如510300买入持有累计收益75.01\%，高于同期沪深300价格指数的54.14\%。该差异包含分红、ETF价格与跟踪表现等因素，不能全部归因于分红。各单ETF账户使用完整的单资产仓位，组合则按$1/K_t$上限缩放，并随新ETF纳入而调整资金分配；独立账户净值的简单平均不等于组合净值。六只ETF的完整净值图列于附录\ref{app:etffig}。

\section{六宽基组合：主要差距来自下跌期与反弹期}
\subsection{整体收益与风险}
固定六ETF组合累计收益41.85\%，期末净值1.418455，年化收益5.52\%，最大回撤24.09\%，年化波动10.67\%。组合在2020年获得较好增长，2022年明显回撤，2023年至2024年前期继续承压，随后随市场回升恢复。分散到多个宽基ETF没有消除同类权益资产共同下跌的风险。
\begin{table}[htbp]\centering
\caption{六宽基ETF组合与原研报结果比较；差值为本研究减研报}
\begin{tabular}{lrrr}\toprule
指标 & 本研究组合 & 研报宽基组合 & 差值（百分点）\\\midrule
年化收益 & 5.52\% & 11.78\% & -6.26\\
最大回撤 & 24.09\% & 11.50\% & +12.59\\
年化波动 & 10.67\% & 8.68\% & +1.99\\\bottomrule
\end{tabular}
\end{table}

相较研报，本研究组合年化收益低约6.26个百分点，最大回撤约为其两倍，波动也更高。这并非单纯因为降低风险而牺牲收益，而是尚未实现研报所报告的收益与回撤组合。需要强调，研报另有全市场动态ETF池策略，年化收益为12.57\%，与这里的六宽基组合属于不同投资范围；本文以11.78\%这一宽基结果为比较对象。
\snapshotfigure{etf_portfolio_nav.png}{固定六宽基ETF组合净值，2019-01-02至2025-09-30。}{fig:portfolio}

\subsection{年度收益揭示了怎样的差距}
\begin{table}[htbp]\centering
\caption{年度收益比较；2025年为1--9月，差值为本研究减研报}
\begin{tabular}{lrrr}\toprule
年份 & 本研究组合 & 研报宽基组合 & 差值（百分点）\\\midrule
2019 & 13.22\% & 21.75\% & -8.53\\
2020 & 18.16\% & 10.56\% & +7.60\\
2021 & 5.47\% & 7.07\% & -1.60\\
2022 & -13.31\% & 5.11\% & -18.42\\
2023 & -2.20\% & 3.06\% & -5.26\\
2024 & 4.79\% & 18.75\% & -13.96\\
2025（1--9月） & 13.14\% & 13.63\% & -0.49\\\bottomrule
\end{tabular}
\end{table}

2022年是最突出的分歧：本研究亏损13.31\%，研报盈利5.11\%，差距约18.42个百分点。2023年本研究继续小幅亏损，说明降低仓位和多宽基分配并不足以在这一阶段形成正收益。网格在价格持续下行时仍可能因逆向加仓承担损失，而策略2仅将中枢从50\%调至40\%，并未切断权益风险敞口；这是与亏损结果相容的机制解释，但不能替代逐日持仓的损益归因。

2024年本研究收益4.79\%，显著低于研报的18.75\%。结合净值曲线的修复过程，收益缺口也涉及上涨参与，而不仅是下跌防守。逆向减仓和两交易日执行滞后可能影响急涨阶段的参与程度，其影响方向与大小仍取决于当时持仓，不能仅凭年度结果量化。2025年前三季度两者收益接近，也不足以证明此前差距已经消失。

2020年本研究收益高于研报，说明偏差并非各年以相同比例缩放。更值得研究的是策略对不同市场路径的响应，尤其是连续下跌中的仓位累积、反弹初期的风险暴露，以及现金与权益之间的切换。年度收益的对照支持“当前实现未获得研报所描述的年度稳定性”，但不能单独识别哪一项规则造成差距。

\subsection{权益仓位与交易摩擦}
组合权益仓位由六个指数信号共同决定，2025年9月30日的计收益仓位为43.33\%。未使用的资金收益设为0，因此组合同时承担权益参与不足和现金收益缺失的影响。原研报将闲置资金配置短融ETF，该部分在本研究中没有收益贡献。
\snapshotfigure{etf_portfolio_equity_weight.png}{六宽基ETF组合权益仓位：各ETF目标权重滞后两交易日计收益。}{fig:allocation}

全样本名义单边权重变化累计约82倍组合资产，显示策略需要较频繁地调整仓位。这里的统计反映目标权重变化，未模拟价格漂移后的真实成交额，不能直接视作实际账户换手。由于费率和滑点均设为0，5.52\%的年化收益尚未扣除交易摩擦；有摩擦条件下的表现还需要单独评估。

组合仅使用策略2，而它并非多数独立ETF的收益最优规则。这个选择有助于展示温和中枢调整的效果，但缺少其他组合规则及同资产池买入持有基准的对照，无法回答策略2在组合层面是否更有效，也无法量化网格相对被动配置的增益。

\section{与研报的差异及结果解释边界}
本研究复用了研报的ATR偏离度、基础网格和四类择时思路，但投资池、策略组合与交易假设并未完全一致。下表列出影响结果解释的主要差异。
\begin{table}[htbp]\centering\small
\caption{主要方法差异及其潜在影响}
\begin{tabularx}{\linewidth}{l Y Y}\toprule
维度 & 本研究设定 & 与研报比较时的含义\\\midrule
资产选择 & 固定六只ETF，上市且信号有效后纳入 & 未还原按历史规模、成交额及同指数最大规模ETF选择的动态过程\\
择时规则 & 四种规则分别测试，组合仅用策略2 & 研报还讨论多规则联合版本，单规则与联合规则不能直接等同\\
指标参数 & ATR与中枢20日、趋势60日 & 研报以$N,N_1$描述窗口，相关页未完整披露数值，参数差异可能改变信号\\
现金配置 & 闲置资金收益0 & 缺少研报短融ETF的持有收益\\
交易执行 & 收盘信号次日收盘执行，两日后计收益 & 研报未明确披露完全相同的执行时点，快速行情中的参与可能不同\\
收益数据 & 本地指数与ETF行情、分红及折算调整 & 与Wind数据、复权及历史事件口径可能不同\\
实务约束 & 费用0、缺失ETF收益填0 & 未完整反映交易成本、停牌调仓及可成交性\\\bottomrule
\end{tabularx}
\end{table}

研报的通用ETF组合流程按月更新备选池，以基金规模20亿元、历史60日平均成交额2亿元为筛选阈值；相同跟踪指数保留规模较大的ETF，剩余资金配短融ETF。宽基版本将范围限制在本文六个指数，但没有列出完整的历史ETF选择及纳入细节。因此，本研究固定ETF池能够研究仓位管理思路，却不足以完整代表研报历史组合。

指数层面也存在比较范围差异。研报统计28只指数的平均累计收益差和回撤改善，并展示联合择时策略；本文为六宽基的独立规则比较。尤其是研报2022年以来的创业板择时曲线期末高于指数并取得正收益，而本文四种独立规则均未达到这一结果。该差异说明还有影响收益路径的设定尚未对齐，不能将机制相近理解为业绩已经复现。

现有比较无法分解每项差异的贡献。后续研究应在同一资产池和交易口径下，分别比较固定仓位、基础网格、独立择时及联合择时，再逐项加入现金收益和交易成本。这样的对照才能区分仓位降低、趋势约束、资产选择与现金管理各自带来的作用。

\section{结论}
ATR将价格偏离转化为随波动调整的仓位规则，为宽基ETF提供了一种明确的风险预算方法。但2019年以来的结果表明，较低回撤主要伴随着较低的上涨参与，基础网格并未在本组长期样本中获得更高的绝对收益，也不能仅以回撤下降认定其优于被动持有。

择时对这一取舍有改善空间。沪深300采用策略3后，年化收益从2.89\%升至5.00\%，夏普从0.3449升至0.5537；然而，创业板、中证1000和科创50的结果说明，收益改善、回撤改善及规则排序并不具有普遍一致性。趋势约束应被理解为需要按资产和市场路径评估的配置机制，而非保证网格盈利的条件。

固定六ETF组合年化收益5.52\%、最大回撤24.09\%，且在2022、2023年连续亏损，尚未重现研报的风险收益表现。当前最清晰的研究发现是：单靠波动自适应网格、温和中枢择时及宽基分散，仍无法同时解决持续下跌中的风险累积与持续上涨中的收益让渡。策略的进一步评估需要围绕风险暴露、联合规则、现金收益和实际交易条件展开。

\appendix
\clearpage
\section{策略与代码对应关系}\label{app:code}
指数指标、基础网格及四种择时由程序\cite{nb1}实现；ETF总回报、单ETF账户及组合由程序\cite{nb2}实现。下列代码索引用于定位策略定义，正文中的策略名称与相应仓位变量一一对应。
\begin{table}[htbp]\centering\small
\caption{指数策略代码索引}
\begin{tabularx}{\linewidth}{l l Y}\toprule
研究环节 & 函数或仓位变量 & 实现位置及内容\\\midrule
ATR与NATR & \code{add_atr} & 指数程序第13步：计算TR、ATR、20日中枢及NATR\\
基础网格 & \code{grid_weight}、\code{w_base} & 第15、17步：将NATR映射至0--100\%目标仓位\\
趋势信号 & \code{add_timing_ma}、\code{timing} & 第24--25步：60日趋势均线和多空方向\\
策略1 & \code{w_t1} & 第26步：按多空信号改变上下侧间距\\
策略2 & \code{w_t2} & 第27步：将仓位中枢设为60\%、40\%或50\%\\
策略3 & \code{w_t3} & 第28步：按信号约束50\%仓位下限或上限\\
策略4 & \code{limit_vs_last}、\code{w_t4} & 第29--30步：递推约束仓位增减方向\\
收益与绩效 & \code{index_backtest}、\code{performance} & 第19、22步：滞后执行收益与绩效统计\\\bottomrule
\end{tabularx}
\end{table}
\begin{table}[htbp]\centering\small
\caption{ETF与组合代码索引}
\begin{tabularx}{\linewidth}{l l Y}\toprule
研究环节 & 函数或变量 & 实现位置及内容\\\midrule
现金分红 & \code{read_dividend} & ETF程序第8步：累计分红差分还原单次事件\\
份额与总回报 & \code{share_ratio}、\code{total_return} & 第13--14步：按份额换算系数修正ETF收益\\
组合纳入 & \code{eligible}、\code{n_eligible} & 第20--21步：上市及信号有效，确定$K_t$\\
目标与持仓 & \code{target}、\code{holdings} & 第22--23步：$w^{(2)}/K_t$及两交易日收益滞后\\
组合收益 & \code{portfolio_return} & 第25步：ETF总回报按持仓加权，扣除设定费用\\
单ETF策略 & \code{daily}、\code{nav} & 第33步：五种仓位分别应用于单ETF总回报\\\bottomrule
\end{tabularx}
\end{table}

指数完整绩效见\path{outputs/features/index_metrics_summary.csv}，单ETF净值见同目录\code{etf_*_individual_nav.csv}，组合逐日结果、年度收益和总体指标分别见\code{portfolio_daily.csv}、\code{portfolio_annual_returns.csv}和\code{portfolio_vs_report.csv}。数据和结果路径用于研究追溯，不改变正文的收益与风险定义。

\clearpage
\section{六指数完整绩效}
下表列示六个指数在两个区间的全部结果。起始年份2019对应2019-01-02，2022对应2022-01-04，截止日均为2025-09-30；科创50的2019栏实际从2020-05-08起。所有策略采用20日ATR与价格中枢、60日趋势判断及零交易费率。收益、回撤和波动保留两位百分数，夏普保留四位小数。
\begingroup\small
\begin{longtable}{ll l rrrr}
\caption{六指数完整绩效：夏普的无风险收益取0}\\\toprule
代码 & 起始年份 & 方法 & 年化收益 & 最大回撤 & 年化波动 & 夏普\\\midrule
\endfirsthead
\toprule 代码 & 起始年份 & 方法 & 年化收益 & 最大回撤 & 年化波动 & 夏普\\\midrule
\endhead
\bottomrule\endfoot
000016 & 2019 & 基础 & 2.48\% & 23.01\% & 9.37\% & 0.3080\\
000016 & 2019 & 择时1 & 3.06\% & 22.51\% & 9.37\% & 0.3685\\
000016 & 2019 & 择时2 & 3.24\% & 21.64\% & 9.30\% & 0.3894\\
000016 & 2019 & 择时3 & 3.44\% & 21.85\% & 9.43\% & 0.4055\\
000016 & 2019 & 择时4 & 3.46\% & 22.35\% & 9.41\% & 0.4081\\
000016 & 2019 & 指数 & 4.16\% & 45.41\% & 18.68\% & 0.3114\\
000016 & 2022 & 基础 & -1.17\% & 17.55\% & 8.96\% & -0.0870\\
000016 & 2022 & 择时1 & -1.24\% & 17.66\% & 8.74\% & -0.0995\\
000016 & 2022 & 择时2 & -0.52\% & 16.66\% & 8.69\% & -0.0166\\
000016 & 2022 & 择时3 & -0.91\% & 15.74\% & 8.60\% & -0.0630\\
000016 & 2022 & 择时4 & -0.88\% & 16.22\% & 8.56\% & -0.0604\\
000016 & 2022 & 指数 & -2.50\% & 32.83\% & 17.33\% & -0.0595\\\midrule
000300 & 2019 & 基础 & 2.89\% & 23.24\% & 9.58\% & 0.3449\\
000300 & 2019 & 择时1 & 4.50\% & 20.41\% & 9.57\% & 0.5080\\
000300 & 2019 & 择时2 & 3.69\% & 21.54\% & 9.55\% & 0.4270\\
000300 & 2019 & 择时3 & 5.00\% & 20.80\% & 9.66\% & 0.5537\\
000300 & 2019 & 择时4 & 4.61\% & 21.28\% & 9.64\% & 0.5162\\
000300 & 2019 & 指数 & 6.88\% & 45.60\% & 19.22\% & 0.4424\\
000300 & 2022 & 基础 & -1.87\% & 20.13\% & 9.21\% & -0.1585\\
000300 & 2022 & 择时1 & -0.96\% & 17.42\% & 8.99\% & -0.0625\\
000300 & 2022 & 择时2 & -0.87\% & 17.14\% & 8.99\% & -0.0524\\
000300 & 2022 & 择时3 & -0.85\% & 16.88\% & 8.89\% & -0.0513\\
000300 & 2022 & 择时4 & -0.84\% & 16.07\% & 8.91\% & -0.0499\\
000300 & 2022 & 指数 & -1.72\% & 36.05\% & 18.14\% & -0.0052\\\midrule
000905 & 2019 & 基础 & 5.02\% & 22.26\% & 11.55\% & 0.4822\\
000905 & 2019 & 择时1 & 5.12\% & 21.59\% & 11.54\% & 0.4904\\
000905 & 2019 & 择时2 & 4.96\% & 21.30\% & 11.54\% & 0.4775\\
000905 & 2019 & 择时3 & 5.67\% & 20.68\% & 11.45\% & 0.5388\\
000905 & 2019 & 择时4 & 5.06\% & 19.66\% & 11.30\% & 0.4933\\
000905 & 2019 & 指数 & 9.26\% & 41.81\% & 21.93\% & 0.5135\\
000905 & 2022 & 基础 & 0.96\% & 21.54\% & 12.06\% & 0.1395\\
000905 & 2022 & 择时1 & 0.61\% & 21.12\% & 11.98\% & 0.1107\\
000905 & 2022 & 择时2 & 0.83\% & 20.99\% & 12.06\% & 0.1287\\
000905 & 2022 & 择时3 & 1.32\% & 19.84\% & 11.60\% & 0.1711\\
000905 & 2022 & 择时4 & 1.42\% & 18.68\% & 11.61\% & 0.1791\\
000905 & 2022 & 指数 & 0.20\% & 39.52\% & 21.93\% & 0.1185\\\midrule
000852 & 2019 & 基础 & 4.54\% & 30.88\% & 13.46\% & 0.3976\\
000852 & 2019 & 择时1 & 4.20\% & 31.37\% & 13.74\% & 0.3686\\
000852 & 2019 & 择时2 & 4.08\% & 31.60\% & 13.54\% & 0.3631\\
000852 & 2019 & 择时3 & 5.37\% & 29.51\% & 13.48\% & 0.4556\\
000852 & 2019 & 择时4 & 4.08\% & 31.88\% & 13.39\% & 0.3656\\
000852 & 2019 & 指数 & 8.60\% & 46.71\% & 24.80\% & 0.4570\\
000852 & 2022 & 基础 & -1.65\% & 30.88\% & 14.55\% & -0.0415\\
000852 & 2022 & 择时1 & -2.57\% & 31.37\% & 14.92\% & -0.0993\\
000852 & 2022 & 择时2 & -2.14\% & 31.54\% & 14.67\% & -0.0739\\
000852 & 2022 & 择时3 & -0.38\% & 29.40\% & 14.30\% & 0.0453\\
000852 & 2022 & 择时4 & -2.19\% & 31.88\% & 14.50\% & -0.0798\\
000852 & 2022 & 指数 & -1.54\% & 46.41\% & 25.36\% & 0.0659\\\midrule
399006 & 2019 & 基础 & 4.92\% & 34.50\% & 14.05\% & 0.4124\\
399006 & 2019 & 择时1 & 5.95\% & 35.61\% & 13.82\% & 0.4870\\
399006 & 2019 & 择时2 & 5.10\% & 35.48\% & 13.81\% & 0.4293\\
399006 & 2019 & 择时3 & 7.70\% & 34.27\% & 13.89\% & 0.6032\\
399006 & 2019 & 择时4 & 6.21\% & 35.29\% & 13.78\% & 0.5063\\
399006 & 2019 & 指数 & 15.75\% & 57.05\% & 29.20\% & 0.6461\\
399006 & 2022 & 基础 & -3.46\% & 32.45\% & 14.26\% & -0.1755\\
399006 & 2022 & 择时1 & -3.14\% & 34.00\% & 13.75\% & -0.1636\\
399006 & 2022 & 择时2 & -3.36\% & 32.84\% & 13.72\% & -0.1807\\
399006 & 2022 & 择时3 & -2.29\% & 32.27\% & 13.56\% & -0.1029\\
399006 & 2022 & 择时4 & -3.25\% & 32.91\% & 13.51\% & -0.1770\\
399006 & 2022 & 指数 & -0.71\% & 53.94\% & 29.82\% & 0.1228\\\midrule
000688 & 2019 & 基础 & 1.77\% & 36.43\% & 14.61\% & 0.1930\\
000688 & 2019 & 择时1 & 3.19\% & 36.51\% & 14.62\% & 0.2880\\
000688 & 2019 & 择时2 & 2.61\% & 35.11\% & 14.45\% & 0.2501\\
000688 & 2019 & 择时3 & 3.24\% & 37.04\% & 14.46\% & 0.2925\\
000688 & 2019 & 择时4 & 2.30\% & 37.87\% & 14.38\% & 0.2296\\
000688 & 2019 & 指数 & 3.91\% & 62.66\% & 30.15\% & 0.2755\\
000688 & 2022 & 基础 & -0.36\% & 32.53\% & 14.92\% & 0.0501\\
000688 & 2022 & 择时1 & -0.61\% & 32.64\% & 14.82\% & 0.0330\\
000688 & 2022 & 择时2 & 0.29\% & 30.54\% & 14.71\% & 0.0931\\
000688 & 2022 & 择时3 & -0.69\% & 32.71\% & 14.40\% & 0.0237\\
000688 & 2022 & 择时4 & -1.12\% & 33.85\% & 14.53\% & -0.0054\\
000688 & 2022 & 指数 & 1.88\% & 54.01\% & 30.40\% & 0.2095\\
\end{longtable}
\endgroup

\clearpage
\section{单ETF策略净值图}\label{app:etffig}
各图将基础网格、四种择时规则与ETF总回报买入持有基准比较，净值均从1开始。前五只ETF基点为2018-12-28，588000为2020-11-16，截止日均为2025-09-30。不同起点的累计收益不宜直接横向排序。
\pairfigure{etf_510050_individual_nav.png}{etf_510300_individual_nav.png}{单ETF策略净值（一）}{510050}{510300}
\pairfigure{etf_510500_individual_nav.png}{etf_512100_individual_nav.png}{单ETF策略净值（二）}{510500}{512100}
\pairfigure{etf_159915_individual_nav.png}{etf_588000_individual_nav.png}{单ETF策略净值（三）}{159915}{588000}

\clearpage
\begin{thebibliography}{9}
\bibitem{gfreport} 广发证券金融工程研究团队. 量化·精选·致远：2026深度研究三——基于平均真实波幅（ATR）的ETF网格交易策略. 指标与指数策略见第6--10页，宽基ETF组合见第17--18页及表14、15.
\end{thebibliography}
\end{document}
